\sectionkicker{Source-backed technical notes}
\subsection*{Inference \& Serving — cache・MoE・kernel・runtimeの最適化: Technical Notes}
本欄は記事本文で圧縮した一次資料上の情報を、比較・再検証しやすい形で整理したものである。時系列、確認済みの主張、留意点、一次資料URLを掲載する。
\medskip
\subsection*{Theme at a glance}
\addcontentsline{toc}{subsection}{Theme at a glance}
\begin{center}
\footnotesize
\begin{tabularx}{\linewidth}{@{}p{0.20\linewidth}p{0.10\linewidth}p{0.14\linewidth}X@{}}
\toprule
資料 & 位置づけ & 種別 & 時系列 \\
\midrule
vllm-project/vllm May 2026 release series & 主要資料 & FRAMEWORK & 2026-05-04 (FRAMEWORK\_RELEASE); 2026-05-15 (FRAMEWORK\_RELEASE); 2026-05-29 (FRAMEWORK\_RELEASE) \\
sgl-project/sglang May 2026 release series & 主要資料 & FRAMEWORK & 2026-05-05 (FRAMEWORK\_RELEASE); 2026-05-16 (FRAMEWORK\_RELEASE) \\
Release v0.6.12 & 主要資料 & FRAMEWORK & 2026-05-29 (FRAMEWORK\_RELEASE) \\
Adaptive KV Cache Reuse for Fast Long-Context LLM Serving & 補足資料 & 論文 & 2026-05-20 (論文\_RELEASE) \\
ReMoE: Boosting Expert Reuse through Router Fine-Tuning in Memory-Constrained MoE LLM Inference & 補足資料 & 論文 & 2026-05-26 (論文\_RELEASE) \\
\bottomrule
\end{tabularx}
\end{center}
\normalsize

\begin{technicalnote}{vllm-project/vllm May 2026 release series}{主要資料}
\begin{tabularx}{\linewidth}{@{}>{\bfseries}p{0.22\linewidth}X@{}}
組織 & vllm-project/vllm \\
種別 & FRAMEWORK \\
時系列 & 2026-05-04 (FRAMEWORK\_RELEASE); 2026-05-15 (FRAMEWORK\_RELEASE); 2026-05-29 (FRAMEWORK\_RELEASE) \\
\end{tabularx}
\smallskip
{\bfseries 一次資料から整理したtechnical points}
\begin{itemize}[leftmargin=1.5em,itemsep=0.35em]
\item \textbf{一次情報で確認できる事実}: The preserved primary source identifies “v0.20.1” on 2026-05-04.
\item \textbf{一次情報で確認できる事実}: The preserved primary source identifies “v0.21.0” on 2026-05-15.
\item \textbf{一次情報で確認できる事実}: The preserved primary source identifies “v0.22.0” on 2026-05-29.
\item \textbf{Project claim}: vLLM's May releases v0.20.1, v0.21.0, and v0.22.0 form a coherent serving series centered on DeepSeek V4 hardening, KV/offload evolution, attention backends, Model Runner V2, and speculative/disaggregated serving.
\end{itemize}
{\bfseries 読む際の境界}
\begin{itemize}[leftmargin=1.5em,itemsep=0.35em]
\item \textbf{分析上の留意点}: Official repository release notes verify release chronology and listed implementation changes; project performance claims and cross-hardware generalization were not independently reproduced.
\end{itemize}
{\bfseries 一次資料}
\begingroup\sloppy
\begin{itemize}[leftmargin=1.5em,itemsep=0.25em]
\item {\scriptsize\url{https://github.com/vllm-project/vllm/releases/tag/v0.20.1}}
\item {\scriptsize\url{https://github.com/vllm-project/vllm/releases/tag/v0.21.0}}
\item {\scriptsize\url{https://github.com/vllm-project/vllm/releases/tag/v0.22.0}}
\end{itemize}
\endgroup
\end{technicalnote}
\medskip

\begin{technicalnote}{sgl-project/sglang May 2026 release series}{主要資料}
\begin{tabularx}{\linewidth}{@{}>{\bfseries}p{0.22\linewidth}X@{}}
組織 & sgl-project/sglang \\
種別 & FRAMEWORK \\
時系列 & 2026-05-05 (FRAMEWORK\_RELEASE); 2026-05-16 (FRAMEWORK\_RELEASE) \\
\end{tabularx}
\smallskip
{\bfseries 一次資料から整理したtechnical points}
\begin{itemize}[leftmargin=1.5em,itemsep=0.35em]
\item \textbf{一次情報で確認できる事実}: The preserved primary source identifies “v0.5.11” on 2026-05-05.
\item \textbf{一次情報で確認できる事実}: The preserved primary source identifies “v0.5.12” on 2026-05-16.
\item \textbf{Project claim}: SGLang's May releases v0.5.11 and v0.5.12 form a coherent serving series covering speculative decoding, disaggregated prefix caching, broad model support, and a substantial DeepSeek V4 inference path.
\end{itemize}
{\bfseries 読む際の境界}
\begin{itemize}[leftmargin=1.5em,itemsep=0.35em]
\item \textbf{分析上の留意点}: Official repository release notes verify release chronology and listed implementation changes; project performance claims and cross-hardware generalization were not independently reproduced.
\end{itemize}
{\bfseries 一次資料}
\begingroup\sloppy
\begin{itemize}[leftmargin=1.5em,itemsep=0.25em]
\item {\scriptsize\url{https://github.com/sgl-project/sglang/releases/tag/v0.5.11}}
\item {\scriptsize\url{https://github.com/sgl-project/sglang/releases/tag/v0.5.12}}
\end{itemize}
\endgroup
\end{technicalnote}
\medskip

\begin{technicalnote}{Release v0.6.12}{主要資料}
\begin{tabularx}{\linewidth}{@{}>{\bfseries}p{0.22\linewidth}X@{}}
組織 & flashinfer-ai/flashinfer \\
種別 & FRAMEWORK \\
時系列 & 2026-05-29 (FRAMEWORK\_RELEASE) \\
\end{tabularx}
\smallskip
{\bfseries 一次資料から整理したtechnical points}
\begin{itemize}[leftmargin=1.5em,itemsep=0.35em]
\item \textbf{一次情報で確認できる事実}: The preserved primary source identifies “Release v0.6.12” on 2026-05-29.
\item \textbf{Project claim}: FlashInfer v0.6.12 adds inference-kernel work including NVFP4 quantization, MoE kernels, GQA/MLA paths, and model/backend support relevant to low-level serving.
\end{itemize}
{\bfseries 読む際の境界}
\begin{itemize}[leftmargin=1.5em,itemsep=0.35em]
\item \textbf{分析上の留意点}: Official repository release notes verify release chronology and listed implementation changes; project performance claims and cross-hardware generalization were not independently reproduced.
\end{itemize}
{\bfseries 一次資料}
\begingroup\sloppy
\begin{itemize}[leftmargin=1.5em,itemsep=0.25em]
\item {\scriptsize\url{https://github.com/flashinfer-ai/flashinfer/releases/tag/v0.6.12}}
\end{itemize}
\endgroup
\end{technicalnote}
\medskip

\begin{technicalnote}{Adaptive KV Cache Reuse for Fast Long-Context LLM Serving}{補足資料}
\begin{tabularx}{\linewidth}{@{}>{\bfseries}p{0.22\linewidth}X@{}}
組織 & - \\
種別 & 論文 \\
時系列 & 2026-05-20 (論文\_RELEASE) \\
\end{tabularx}
\smallskip
{\bfseries 一次資料から整理したtechnical points}
\begin{itemize}[leftmargin=1.5em,itemsep=0.35em]
\item \textbf{一次情報で確認できる事実}: The preserved primary source identifies “Adaptive KV Cache Reuse for Fast Long-Context LLM Serving” on 2026-05-20.
\item \textbf{Author claim}: CacheTune extends KV-cache reuse beyond strict prefixes by adaptively recomputing context needed to preserve attention relationships in long-context serving.
\end{itemize}
{\bfseries 読む際の境界}
\begin{itemize}[leftmargin=1.5em,itemsep=0.35em]
\item \textbf{分析上の留意点}: The preserved original arXiv metadata/abstract is primary-paper evidence, but experimental results and generalization are author-reported and were not independently reproduced in this Evidence review.
\end{itemize}
{\bfseries 一次資料}
\begingroup\sloppy
\begin{itemize}[leftmargin=1.5em,itemsep=0.25em]
\item {\scriptsize\url{http://arxiv.org/abs/2605.24022v1}}
\end{itemize}
\endgroup
\end{technicalnote}
\medskip

\begin{technicalnote}{ReMoE: Boosting Expert Reuse through Router Fine-Tuning in Memory-Constrained MoE LLM Inference}{補足資料}
\begin{tabularx}{\linewidth}{@{}>{\bfseries}p{0.22\linewidth}X@{}}
組織 & - \\
種別 & 論文 \\
時系列 & 2026-05-26 (論文\_RELEASE) \\
\end{tabularx}
\smallskip
{\bfseries 一次資料から整理したtechnical points}
\begin{itemize}[leftmargin=1.5em,itemsep=0.35em]
\item \textbf{一次情報で確認できる事実}: The preserved primary source identifies “ReMoE: Boosting Expert Reuse through Router Fine-Tuning in Memory-Constrained MoE LLM Inference” on 2026-05-26.
\item \textbf{Author claim}: ReMoE fine-tunes MoE routing to increase expert reuse when only a subset of experts can remain cached under memory-constrained inference.
\end{itemize}
{\bfseries 読む際の境界}
\begin{itemize}[leftmargin=1.5em,itemsep=0.35em]
\item \textbf{分析上の留意点}: The preserved original arXiv metadata/abstract is primary-paper evidence, but experimental results and generalization are author-reported and were not independently reproduced in this Evidence review.
\end{itemize}
{\bfseries 一次資料}
\begingroup\sloppy
\begin{itemize}[leftmargin=1.5em,itemsep=0.25em]
\item {\scriptsize\url{http://arxiv.org/abs/2605.27081v1}}
\end{itemize}
\endgroup
\end{technicalnote}
\medskip
